\section{Statement of results}\label{section2}

Our conventions for the paper are as follows.
An algebra is meant to be associative and have a~1.
A subalgebra has the same~1 as the parent algebra.
We f\/ix a f\/ield~$\mathbb F$. All unadorned tensor products
are meant to be over $\mathbb F$. We f\/ix $q \in \mathbb F$ such that
$q^4\not=1$. Recall the natural numbers
$\mathbb N = \lbrace 0,1,2,\ldots \rbrace$ and
integers $\mathbb Z=\lbrace 0, \pm 1, \pm 2,\ldots \rbrace$.

\begin{Definition}[\protect{\cite[Def\/inition~1.2]{uaw}}]
\label{def:uaw}
Def\/ine an $\mathbb F$-algebra
$\Delta=\Delta_q$
by generators and relations in the following way.
The generators are $A$, $B$, $C$.
The relations assert that each of
\begin{gather}
 A+ \frac{qBC-q^{-1}CB}{q^2-q^{-2}},
\qquad B+
\frac{qCA-q^{-1}AC}{q^2-q^{-2}},
\qquad
C+
\frac{qAB-q^{-1}BA}{q^2-q^{-2}}
\label{eq:comlist}
\end{gather}
is central in $\Delta$.
The algebra
$\Delta$ is called the {\it universal Askey--Wilson algebra}.
\end{Definition}

\begin{Definition}[\protect{\cite[Def\/inition~1.3]{uaw}}]
\label{def:abc}
For the
three central elements in
(\ref{eq:comlist}), multiply each by $q+q^{-1}$ to get
$\alpha$, $\beta$, $\gamma$. Thus
\begin{gather}
A+ \frac{qBC-q^{-1}CB}{q^2-q^{-2}}
 =
\frac{\alpha}{q+q^{-1}},
\label{eq:u1}
\\
B+
\frac{qCA-q^{-1}AC}{q^2-q^{-2}}
 =
\frac{\beta}{q+q^{-1}},
\label{eq:u2}
\\
C+
\frac{qAB-q^{-1}BA}{q^2-q^{-2}}
 =
\frac{\gamma}{q+q^{-1}}.
\label{eq:u3}
\end{gather}
Note that each of $\alpha$, $\beta$, $\gamma$ is central in~$\Delta$.
\end{Definition}

We mention a few facts about~$\Delta$.
Recall that the modular group
${\rm  {PSL}}_2(\mathbb Z)$ has a presentation
by generators $\rho$, $\sigma$ and relations
$\rho^3=1$, $\sigma^2=1$. See for example
\cite{alpern}.
By \cite[Theorem~3.1]{uaw},
the group
${\rm  {PSL}}_2(\mathbb Z)$ acts on
$\Delta$ as a group of automorphisms
such that
$\rho$
 sends
$(A,B,C)\mapsto (B,C,A)$
and  $\sigma $
sends
$(A,B,\gamma)\mapsto (B,A,\gamma)$.
By
\cite[Theorem~3.13]{uaw}
this action is faithful.

By \cite[Theorem~4.1]{uaw} the following
is a basis for the $\mathbb F$-vector space $\Delta$:
\begin{gather*}
 A^iB^jC^k \alpha^r\beta^s\gamma^t,
\qquad
i,j,k,r,s,t \in \mathbb N.
\end{gather*}
There is a related basis
\cite[Theorem~7.5]{uaw} that we will use
 in
Section~\ref{section9} below. This related basis
involves
 a
central element $\Omega$ known as the Casimir
element~\cite[Lemma~6.1]{uaw}.
This element is def\/ined
as follows.


\begin{Definition}[\protect{\cite[Lemma~6.1]{uaw}}]
\label{def:casdelta}
Def\/ine $\Omega \in \Delta$ by
\begin{gather*}
\Omega = qABC+q^2A^2+q^{-2}B^2+q^2 C^2 - q A\alpha  - q^{-1} B \beta
- q C \gamma.
\end{gather*}
We call $\Omega$ the {\it Casimir element} of~$\Delta$.
\end{Definition}

\begin{Lemma}[\protect{\cite[Theorem~6.2, Corollary~8.3]{uaw}}]
The Casimir element $\Omega$ is contained in the
center~$Z(\Delta)$.
Moreover
$\lbrace \Omega^i\alpha^r\beta^s\gamma^t \,|\, i,r,s,t \in \mathbb N\rbrace$
is a basis for
the $\mathbb F$-vector space
$Z(\Delta)$,
provided that~$q$ is not a root of unity.
\end{Lemma}

\begin{Lemma}[\protect{\cite[Theorem~6.4]{uaw}}]
\label{lem:psl}
The Casimir element $\Omega$  is fixed by
everything in
${\rm  {PSL}}_2(\mathbb Z)$.
\end{Lemma}

We will be discussing how
$\Delta$ is related to the
quantum universal enveloping
algebra $U_q(\mathfrak{sl}_2)$.
For
this algebra
there are two presentations of
 interest to us;
 the Chevalley presentation
\cite[Section~1.1]{jantzen}
and the equitable presentation~\cite{equit}. We now recall the
  Chevalley presentation.


\begin{Definition}[\protect{\cite[Section~1.1]{jantzen}}]
\label{def:chev}
The
$\mathbb F$-algebra $U=U_q(\mathfrak{sl}_2)$
is def\/ined by generators
$e$, $f$, $k^{\pm 1 }$  and relations
\begin{gather*}
  k k^{-1} = k^{-1} k = 1,
\qquad
 ke = q^2 e k,
\qquad
 kf = q^{-2} f k,
\qquad ef-fe = \frac{k-k^{-1}}{q-q^{-1}}.
\end{gather*}
We call  $e$, $f$, $k^{\pm 1}$ the {\it Chevalley generators} for~$U$.
\end{Definition}

 We now brief\/ly discuss some
f\/inite-dimensional $U$-modules. Strictly speaking
we will not use this information; it is included
in order to clarify the nature of the Casimir element
for $U$ described below.

Recall the notation
\begin{gather*}
  \lbrack  n \rbrack_q=\frac{q^n-q^{-n}}{q-q^{-1}},
  \qquad n \in \mathbb N.
\end{gather*}

\begin{Lemma}[\protect{\cite[Section~2]{jantzen}}]
For all integers $n\geq 0$ and $\varepsilon \in \lbrace 1,-1\rbrace$
there exists a
$U$-module
$L(n,\varepsilon)$ with the following properties.
$L(n,\varepsilon)$ has a basis
$\lbrace v_i\rbrace_{i=0}^n$
such that
\begin{gather*}
k v_i  =  \varepsilon q^{n-2i} v_i, \quad  0 \leq i \leq n,
\\
f v_i = \lbrack i+1 \rbrack_q v_{i+1}, \quad 0 \leq i \leq n-1,
\qquad fv_n =0,
\\
e v_i = \varepsilon
\lbrack n-i+1 \rbrack_q v_{i-1}, \quad 1 \leq i \leq n,
\qquad ev_0 =0.
\end{gather*}
The
$U$-module
$L(n,\varepsilon)$ is irreducible provided that~$q$ is not a root of unity.
\end{Lemma}


In Def\/inition~\ref{def:casdelta} we gave the Casimir element
for $\Delta$. We now recall the Casimir element
for~$U$.

\begin{Definition}[\protect{\cite[Section~2.7]{jantzen}}]
\label{def:casu}
Def\/ine
 $\Phi \in
U$ as follows:
\begin{gather*}
%\label{def:phi}
\Phi = ef +\frac{q^{-1}k + qk^{-1}}{(q-q^{-1})^2}.
\end{gather*}
We call $\Phi$ the {\it Casimir element} of
$U$.
\end{Definition}

\begin{Lemma}[\protect{\cite[Lemma~2.7, Proposition~2.18]{jantzen}}]
The element
$\Phi$ is contained in the center~$Z(U)$.
Moreover $\lbrace \Phi^i\rbrace_{i \in \mathbb N}$ is a basis
for the $\mathbb F$-vector space~$Z(U)$,
 provided that $q$ is not a root of unity.
\end{Lemma}





\begin{Lemma}
\label{lem:normmeaning}
{\rm \cite[Lemma~2.7]{jantzen}.}
On the
$U$-module
$L(n,\varepsilon)$,
\begin{gather*}
\Phi  =  \varepsilon \frac{q^{n+1}+q^{-n-1}}{(q-q^{-1})^2} I.
\end{gather*}
Here $I$ denotes the identity map.
\end{Lemma}

For notational convenience we now adjust the
normalization for $\Phi$.

\begin{Definition}
\label{def:norm}
Def\/ine
\begin{gather}
\Lambda  = \big(q-q^{-1}\big)^2 \Phi
= \big(q-q^{-1}\big)^2ef +q^{-1}k + qk^{-1}.
\label{eq:lambdaform}
\end{gather}
Note that on
$L(n,\varepsilon)$,
\begin{gather*}
\Lambda  =  \varepsilon\big(q^{n+1}+q^{-n-1}\big)I.
%\label{eq:normform}
\end{gather*}
We call
$\Lambda$ the {\it normalized Casimir element} for
$U$.
\end{Definition}


We now recall the equitable presentation for
$U$ \cite{equit}.

\begin{Proposition}[\protect{\cite[Theorem~2.1]{equit}}]
\label{prop:equit}
The algebra
$U$ is isomorphic to the
$\mathbb F$-algebra defined by generators
$x$, $y^{\pm 1}$, $z$ and relations
\begin{gather}
 y y^{-1} = y^{-1} y  = 1,
\label{eq:eq0}
\\
 \frac{qxy-q^{-1}yx}{q-q^{-1}} = 1,
\label{eq:eq1}
\\
 \frac{qyz-q^{-1}zy}{q-q^{-1}} = 1,
\label{eq:eq2}
\\
 \frac{qzx-q^{-1}xz}{q-q^{-1}} = 1.
\label{eq:eq3}
\end{gather}
An isomorphism with the presentation in
Definition~{\rm \ref{def:chev}} is given by
\begin{gather*}
 y^{\pm 1} \mapsto k^{\pm 1},
\qquad
 z \mapsto k^{-1} + f\big(q-q^{-1}\big),
\qquad
 x \mapsto k^{-1} - ek^{-1}q^{-1}\big(q-q^{-1}\big).
\end{gather*}
The inverse of this isomorphism is given by
\begin{gather*}
 k^{\pm 1} \mapsto y^{\pm 1},
\qquad
 f\mapsto (z-y^{-1})\big(q-q^{-1}\big)^{-1},
\qquad
 e\mapsto (1-xy)q\big(q-q^{-1}\big)^{-1}.
\end{gather*}
\end{Proposition}


\begin{Definition}[\protect{\cite[Def\/inition~2.2]{equit}}]
By the {\it equitable presentation} of~$U$ we mean the
presentation given in Proposition
\ref{prop:equit}. We call $x$, $y^{\pm 1}$, $z$ the
{\it equitable generators} for~$U$.
\end{Definition}


\begin{Note}
\label{note:ident}
In what follows we identify the copy of
$U$ given in
Def\/inition~\ref{def:chev}
with
the copy given in
Proposition~\ref{prop:equit},
via the isomorphism given in
Proposition~\ref{prop:equit}.
\end{Note}


In the equitable presentation
of $U$
the normalized Casimir element $\Lambda$ looks as follows.

\begin{Lemma}
\label{lem:sixforms}
The normalized Casimir element $\Lambda$ is equal to each of
the following:
\begin{alignat}{3}
& qx+q^{-1}y+qz-qxyz,
\qquad &&
q^{-1}x+qy+q^{-1}z-q^{-1}zyx, &
\label{eq:v1}
\\
& qy+q^{-1}z+qx-qyzx,
\qquad &&
q^{-1}y+qz+q^{-1}x-q^{-1}xzy,&
\label{eq:v2}
\\
& qz+q^{-1}x+qy-qzxy,
\qquad &&
q^{-1}z+qx+q^{-1}y-q^{-1}yxz.&
\label{eq:v3}
\end{alignat}
\end{Lemma}

\begin{proof}
For the data
(\ref{eq:v1})--(\ref{eq:v3})
let
$\Lambda^-_y$, $\Lambda^-_z$, $\Lambda^-_x$
denote the
expressions in the f\/irst column
and let
$\Lambda^+_y$, $\Lambda^+_z$, $\Lambda^+_x$ denote
the
expressions in the second column.
Consider the expression for
$\Lambda$ given in~(\ref{eq:lambdaform}).
Writing this expression in terms of $x$, $y$, $z$ using
the isomorphism in
Proposition~\ref{prop:equit} and Note~\ref{note:ident},
we obtain
$\Lambda=\Lambda^-_y$.
The element $\Lambda^-_y-\Lambda^+_z$ is equal to
$(q-q^{-1})x$ times
\begin{gather}
1 - \frac{qyz-q^{-1}zy}{q-q^{-1}}.
\label{eq:zer1}
\end{gather}
The expression
(\ref{eq:zer1}) is zero by
(\ref{eq:eq2}) so
$\Lambda^-_y=\Lambda^+_z$.
Similarly one f\/inds
$\Lambda^-_z=\Lambda^+_x$ and
$\Lambda^-_x=\Lambda^+_y$.
The element $\Lambda^-_y-\Lambda^+_x$ is equal to
\begin{gather}
1 - \frac{qxy-q^{-1}yx}{q-q^{-1}}
\label{eq:zer2}
\end{gather}
times
$(q-q^{-1})z$.
The expression~(\ref{eq:zer2}) is zero by~(\ref{eq:eq1}) so
$\Lambda^-_y=\Lambda^+_x$.
Similarly one f\/inds
$\Lambda^-_z=\Lambda^+_y$ and
$\Lambda^-_x=\Lambda^+_z$.
By these comments $\Lambda$ is equal to each of
$\Lambda^{\pm }_x$, $\Lambda^{\pm }_y$, $\Lambda^{\pm }_z$.
\end{proof}



Recall that $a$, $b$, $c$ are mutually commuting indeterminates.
Let
$\mathbb F \lbrack a^{\pm 1}, b^{\pm 1}, c^{\pm 1}\rbrack$
denote the $\mathbb F$-algebra
of Laurent polynomials in
$a$, $b$, $c$ that have all coef\/f\/icients in~$\mathbb F$.

We now state our main results.


\begin{Theorem}\label{thm:main}
There exists a unique  $\mathbb F$-algebra homomorphism
$\natural: \Delta \to
U \otimes
\mathbb F \lbrack a^{\pm 1}, b^{\pm 1}, c^{\pm 1}\rbrack
$ that sends
\begin{gather*}
A  \mapsto   x \otimes a + y \otimes a^{-1} +
\frac{x y - y x}{q-q^{-1}} \otimes b c^{-1},
\\
B \mapsto  y \otimes b + z \otimes b^{-1} +
\frac{y z - z y}{q-q^{-1}} \otimes c a^{-1},
\\
C \mapsto  z \otimes c + x \otimes c^{-1} +
\frac{z x - x z}{q-q^{-1}} \otimes a b^{-1},
\end{gather*}
where $x$, $y$, $z$ denote the equitable generators for~$U$.
The homomorphism $\natural$ sends
\begin{gather}
\alpha  \mapsto  \Lambda \otimes \big(a+a^{-1}\big) + 1\otimes \big(b+b^{-1}\big)\big(c+c^{-1}\big),
\label{eq:alsend}
\\
\beta  \mapsto  \Lambda \otimes \big(b+b^{-1}\big) + 1 \otimes \big(c+c^{-1}\big)\big(a+a^{-1}\big),
\label{eq:besend}
\\
\gamma  \mapsto  \Lambda \otimes \big(c+c^{-1}\big) + 1 \otimes \big(a+a^{-1}\big)\big(b+b^{-1}\big),
\label{eq:gasend}
\end{gather}
where $\Lambda $ denotes the normalized Casimir element of~$U$.
\end{Theorem}


\begin{Theorem}\label{thm:main2}
Under the homomorphism
$\natural$
from Theorem~{\rm \ref{thm:main}}, the image of
$\Omega$ is
\begin{gather}
1 \otimes \big(q+q^{-1}\big)^2
-
1\otimes \big(a+a^{-1}\big)^2
-
1\otimes \big(b+b^{-1}\big)^2
-
1 \otimes \big(c+c^{-1}\big)^2
\nonumber\\
\phantom{1 \otimes \big(q+q^{-1}\big)^2}{}
-\Lambda \otimes \big(a+a^{-1}\big)\big(b+b^{-1}\big)\big(c+c^{-1}\big)
- \Lambda ^2 \otimes 1.\label{eq:omimage}
\end{gather}
Here $\Lambda$ denotes the normalized Casimir element
of~$U$.
\end{Theorem}

\begin{Theorem}
\label{thm:main3}
The
homomorphism~$\natural$
from Theorem~{\rm \ref{thm:main}} is injective.
\end{Theorem}


We mention a corollary to
Theorem~\ref{thm:main3}.
For an $\mathbb F$-algebra $\mathcal A$,
an element $u \in {\mathcal A}$ is called a~{\it zero divisor}
whenever $u\not=0$ and
there exists $0 \not=v \in \mathcal A$
such that~$uv=0$.
By~\cite[Proposition~1.8]{jantzen}
the
algebra $U$
contains no zero divisors.
For an $\mathbb F$-algebra $\mathcal A$
and indeterminate $\lambda$ consider the $\mathbb F$-algebra
${\mathcal A}\otimes \mathbb F \lbrack \lambda,\lambda^{-1}\rbrack$.
One checks that
$\mathcal A$ contains no zero divisors if and only if
${\mathcal A}\otimes \mathbb F \lbrack \lambda,\lambda^{-1}\rbrack$
contains no zero divisors.
Applying this comment three times we see that
the algebra
$
U \otimes
\mathbb F \lbrack a^{\pm 1}, b^{\pm 1}, c^{\pm 1}\rbrack
$ contains no zero divisors.
By this and
Theorem~\ref{thm:main3} we obtain the following result.

\begin{Corollary}
\label{cor:nozerodiv}
The $\mathbb F$-algebra $\Delta$ contains no zero divisors.
\end{Corollary}
